\section{引言：两只小猫的故事}

Jim Fan 是 NVIDIA 的高级研究科学家，领导 Project GR00T（通用机器人基础模型）。他以 1963 年 Held 和 Hein 的经典实验开场：两只新生小猫在旋转木马装置中，主动小猫能自由移动并观察世界，被动小猫只能被动观看——结果只有主动小猫发育出健康的视觉系统。

\begin{importantbox}{核心隐喻：具身交互是智能的基础}
被动观察不足以产生真正的理解——智能需要与物理世界的\textbf{主动交互}。这正是 LLM（被动小猫，只阅读互联网文本）和具身 Agent（主动小猫，与物理世界交互）的根本区别。
\end{importantbox}

\subsection{本章小结}

真正的 AGI 不能仅靠阅读文本实现，需要具身智能（embodied intelligence）——Agent 必须能感知、行动并从物理世界的反馈中学习。

